\documentclass[a4j,11pt,dvipdfmx]{jsarticle}
\usepackage{float}
\usepackage{array}
\usepackage{titlesec}
\usepackage[dvipdfmx]{graphicx}
\usepackage{graphicx}
\usepackage{comment}
\usepackage{multirow}
\usepackage{array}
\usepackage{longtable}
\usepackage{float}

\titleformat*{\section}{\bfseries}
\titleformat*{\subsection}{\bfseries}

\makeatletter
  \renewcommand{\section}{%
    \@startsection{section}{1}{\z@}%
    {0.4\Cvs}{0.1\Cvs}%
    {\normalfont\headfont\raggedright}}
\makeatother

\makeatletter
  % sectionの下マージンを小さく
  \renewcommand{\subsection}{%
    \@startsection{subsection}{1}{\z@}%
    {0.1\Cvs}{0.1\Cvs}%
    {\normalfont\headfont\raggedright}}
\makeatother

\renewcommand{\thesubsection}{\thesection-\arabic{subsection}.}


%---------------------------------------------------
% ページの設定
%---------------------------------------------------
\setlength{\textwidth}{160truemm}
\setlength{\textheight}{250truemm}
\setlength{\topmargin}{-4.5truemm}
\setlength{\oddsidemargin}{0.5truemm}
\pagestyle{empty}
\setlength{\headheight}{0truemm}
\setlength{\parindent}{1zw}
\newcolumntype{b}{!{\vrule width 1pt}}
\newcommand{\bhline}{\noalign{\hrule height 1pt}}   


\begin{document}
提出日: \today
\begin{center}
%\huge{実験計画書}
\huge{第8回ハッカソン議事録}
\end{center}

\begin{table}[htbp]
    \begin{tabular}{llrr}
    \textgt{開催日時}: & \today 12:00〜16:00& &\\
    \textgt{開催場所}: &津田沼校舎 7号館3F 732講義室& &\\

     \textgt{班名}:& 23班 & & \\
     \textgt{メンバー}:& 24G1038 片岡聖 & & \\%k
                      & 24G1075 地曳賢人 (議事録作成者) & & \\%z
                    & 24G1131 御代川稜 & & \\%m
                    & 25G1021 梅澤颯太 & & \\%u
                    & 25G1053 齋藤翔太 & & \\%s
                    & 25G1065 塩澤匠生 & & \\%ta
        \textgt{欠席者}: & なし & &  \\
        \textgt{目的}: &第8週授業までの1週間の各自取り組み内容と進捗確認&&\\
        &および発生した課題の共有と今後の調整の議論．

    \end{tabular}
\end{table}

%%%%%%%%%%%%%%%%%%%%%%%%%%%%%%%%%%%%%%%%%%%%%%%%%
%%%%%%%%%%%%%%%%%%%%%%%%%%%%%%%%%%%%%%%%%%%%%%%%%
%%%%%%%%%%%%%%%%%%%%%%%%%%%%%%%%%%%%%%%%%%%%%%%%%

\section{議題一覧}
 \begin{itemize}
    \item {今週までの成果報告}
    \item {今回の実施内容詳細}
    \item {結合テスト}
    \item{これからの実施内容}
\end{itemize}
\section{討議内容}
\subsection{今週までの成果報告}
\textbf{梅澤}
\begin{itemize}
	\item{作業時間は約1.5〜2時間}
	\item {指揮棒作成のためのモデリング練習}
	\item {性能評価の確認のために，MOPのマークダウン方式の確認}
\end{itemize}

\textbf{塩澤}
\begin{itemize}
	\item{作業時間は約2時間}
	\item{回路図作成}
	\item{輪唱の改善}
\end{itemize}

\textbf{齋藤}
\begin{itemize}
	\item{作業時間は約2時間}
	\item{音色全音の微調整}
	\item{ドラム音改善}
\end{itemize}

\subsection{今回の実施内容詳細}
今回のメインタスクとして，指揮棒回路のはんだ付けを含めた結合を行う\par
梅澤→指揮棒用3Dモデル作成\par
塩澤→はんだ付けを含む実機作成\par
齋藤→別楽器音作成\\

齋藤：現在トランペット，ホルン，トロンボーン，チューバ，ドラムの5楽器だが，全て金管のため音の違いがわかりにくいのではないか．\par
御代川：確かにね，実際のオーケストラではどうなの？\\
　地曵：実際は弦主体で木管金管は規模にもよるけど，各パート2,3人とかチューバは1人ということもある\\
　片岡：じゃあ今余裕あるし，他の楽器の作成も試してみてもいいかもね\\
　齋藤：では弦楽器の作成を試してみます
\subsection{結合テスト}



\section{TAとの議論}
\begin{itemize}
	\item{各メンバーの役割，現状確認}
	\item{Q1. システム要件は達成されているか - 輪唱やドラムのリズム，テンポ変更など}
	\item{A1. 単体テストや仮結合でも問題ないことが確認できている}
	\item{Q2. 予定より早まって進行しているため，要素の追加を行なっているが，最低ラインを設定してるか}
	\item{A2. 音色のクオリティ→ゲーム性とラインはクリアしているが，完全な結合ができていないため，直近で結合していく予定である}
\end{itemize}


\section{次回の会議の日時・場所・予定議題}

\begin{table}[htbp]
    \begin{tabular}{llrr}
    \textgt{開催日時}: & 2026年6月24日12:00〜16:00 & &\\
    \textgt{開催場所}: &津田沼校舎 7号館3F 732講義室& &\\

        \textgt{予定議題}: &実機での本格的な結合テストおよび数値目標達成確認\\
    \end{tabular}
\end{table}
\section{その他，特記事項}
%未定

はんだ付けで問題が発生し，独自に購入したマイコンが破損．\par
既に購入し，来週には手元に届く予定であるが今週はTAへのデモ披露は実施不可．


\end{document}
